\documentclass[9pt,landscape]{article}
\usepackage{{./../static/style}}
\pdfinfo{
  /Title (PyPrimeMesh Cheat Sheet)
  /Creator (TeX)
  /Producer (pdfTeX 1.40.0)
  /Author (Ansys)
  /Subject (PyAnsys)
  /Keywords (PyAnsys, Cheat sheet, template)}

\begin{document}
\raggedright
\footnotesize

Cheat sheet for PyPrimeMesh
% ----------------------------------------

\begin{center}
     \Huge{\textbf{Cheat sheet for PyPrimeMesh}} \\
\end{center}
\begin{center}
  \small{\textbf{version:0.3.0 (stable) }} \\
\end{center}
\AddToShipoutPicture*
  {\put(670,577.5){\includegraphics[height = 1.2cm]{ansys.png}}}
\AddToShipoutPictureBG*{\includegraphics[width=\paperwidth]{bground.png}}
\vspace{-0.15cm}
\noindent\makebox[\linewidth]{\rule{\paperwidth}{2pt}}

\begin{multicols}{3}
\setlength{\premulticols}{1pt}
\setlength{\postmulticols}{1pt}
\setlength{\multicolsep}{1pt}
\setlength{\columnsep}{2pt}

% session starts here. 
% First colomn
% --------------------------------------------------------------------------------
% row 1 col 1
\vfill
\section{\includegraphics[height=\fontcharht\font`\S]{slash.png} Launch PyPrimeMesh Client}
To launch and exit PyPrimeMesh server from Python in gRPC mode.\\
\begin{lstlisting}[language=Python]
# Launch PyPrimeMesh Server, 
from ansys.meshing import prime
with prime.launch_prime(timeout = 20) as prime_client:
  # The Client class gets the model parameter 
  model = prime_client.model
  #Define  Script Here
  # Disconnect from server
  prime_client.exit() 
\end{lstlisting}

To launch an instance of PyPrimeMesh at IP
127.0.0.1 and port 50055 with number of processes 4.
\begin{lstlisting}[language=Python]
 with prime.launch_prime(ip = "127.0.0.1", port = 50055, n_procs = 4) as prime_client:
  model = prime_client.model
\end{lstlisting}

% row 2 col 1
\section{\includegraphics[height=\fontcharht\font`\S]{slash.png} Read and Write files }
Read or Write files of different formats based on file extension.
An example for the same can be seen below\\
\begin{lstlisting}[language=Python]
# Import lucid class
from ansys.meshing.prime import lucid
# Define mesher object
mesher = lucid.Mesh(model)
# Read Mesh (*.msh) file
mesh_file_name = r"sample1_mesh.msh"
mesher.read(mesh_file_name, append=False)
# Import CAD (*.step) file
cad_file_name = r"sample2_CAD.stp"
mesher.read(cad_file_name, append=False)
# Write Mesh (*.cdb) file
cdb_file_name = r"sample3_case.cdb"
  mesher.write(cdb_file_name)
\end{lstlisting}

% row 3 col 1
\section{\includegraphics[height=\fontcharht\font`\S]{slash.png} Defining Size controls }
Sets the global sizing parameters , 

\begin{lstlisting}[language=Python]
# Define Global Size Controls
model.set_global_sizing_params(prime.GlobalSizingParams(min = 0.5, max = 16.0,
                       growth_rate= 1.2))
\end{lstlisting}

Curvature size control definition
\begin{lstlisting}[language=Python]
# Create Curvature Size Control
curvature_control = model.control_data.create_size_control(prime.SizingType.CURVATURE)
control_name = "Curvature_Size_Control"
curvature_control.set_suggested_name(control_name)
# Define Scope 
scope= prime.ScopeDefinition(model, evaluation_type=prime.ScopeEvaluationType.LABELS,
                            label_expression="*")
curvature_control.set_scope(scope)
# Define Curvature Control parameter
curvature_control.set_curvature_sizing_params(prime.CurvatureSizingParams(model, normal_angle=18))
\end{lstlisting}
\section{\includegraphics[height=\fontcharht\font`\S]{slash.png} Generate Wrapper Surface Mesh}
Generate Wrapper Surface Mesh
\begin{lstlisting}[language=Python]
mesher.wrap(
            min_size=0.2, max_size=1.0,
            input_parts="flange,pipe",
            use_existing_features=True,
            recompute_remesh_sizes=True,
            remesh_size_controls=[curvature_control],
          )
\end{lstlisting}

\section{\includegraphics[height=\fontcharht\font`\S]{slash.png} Generate Surface Mesh}
Generate Surface Mesh based on min, max size  
\begin{lstlisting}[language=Python]
surface_scope = prime.ScopeDefinition(model, part_expression="*")
mesher.surface_mesh(min_size=0.5, max_size=16, 
                    generate_quads= True,
                    scope = surface_scope)
\end{lstlisting}

Generate Surface Mesh based on size controls  
\begin{lstlisting}[language=Python]
control_name = "Curvature_Size_Control"
mesher.surface_mesh_with_size_controls(
                                    control_name, 
                                    surface_scope
                                    )
\end{lstlisting}

\section{\includegraphics[height=\fontcharht\font`\S]{slash.png} Surface Mesh Analysis}
Surface Mesh Diagnostics 
\begin{lstlisting}[language=Python]
diag = prime.SurfaceSearch(model)
diag_params = prime.SurfaceDiagnosticSummaryParams(
    model, scope=surface_scope,
    compute_free_edges=True,
) 
diag_res = diag.get_surface_diagnostic_summary(diag_params)
>>> print("Number of free faces : ", diag_res.n_free_faces)
\end{lstlisting}

Surface Mesh Quality Metrics
\begin{lstlisting}[language=Python]
face_quality_measures = prime.FaceQualityMeasure.SKEWNESS
quality = prime.SurfaceSearch(model)
quality_params = prime.SurfaceQualitySummaryParams(
    model=model,
    scope=surface_scope,
    face_quality_measures=[face_quality_measures],
    quality_limit=[0.9],
)
>>> print("Maximum surface skewness : ", qual_summary_res.quality_results[0].max_quality)
\end{lstlisting}

\section{\includegraphics[height=\fontcharht\font`\S]{slash.png} Generate Volume Mesh}
Generate Volume Mesh
\begin{lstlisting}[language=Python]
vol_scope = prime.ScopeDefinition(model, part_expression="*")
mesher.volume_mesh(volume_fill_type=prime.VolumeFillType.HEXCOREPOLY,
                  scope = vol_scope)
\end{lstlisting}
\section{\includegraphics[height=\fontcharht\font`\S]{slash.png} Part Summary}
Query for Part Sumamry
\begin{lstlisting}[language=Python]
part = model.get_part_by_name("sample_part")
summary = part.get_summary(prime.PartSummaryParams(model))
>>> print("Total number of cells: ",summary.n_cells)
\end{lstlisting}

% Add subsection
% This section includes useful links to the documentation.
% Examples: installation, API reference, commands, examples.
% Replace 'name of link' with appropriate display text.

\subsection{References from PyAnsys Documentation}
\begin{itemize}
\item \href{https://prime.docs.pyansys.com/version/stable/getting_started/index.html}{\color{blue}{PyPrimeMesh - Getting Started}}
\item \href{https://prime.docs.pyansys.com/version/stable/user_guide/index.html}{\color{blue}{PyPrimeMesh - User guide}}
\item \href{https://prime.docs.pyansys.com/version/stable/examples/index.html}{\color{blue}{PyPrimeMesh - Examples}}
\end{itemize}
\end{multicols}

% Footer session of the latex with link to documentation and GitHub page
\vspace{-0.15cm}
\noindent\makebox[\linewidth]{\rule{\paperwidth}{4pt}}
\begin{center}
Getting Started with PyAnsys \includegraphics[height=\fontcharht\font`\S]{slash.png} \href{https://github.com/pyansys}{\color{blue}{PyAnsys on GitHub}} \includegraphics[height=\fontcharht\font`\S]{slash.png} Visit \code{\href{https://dev.docs.pyansys.com/}}{\color{blue}{dev.docs.pyansys.com}}
\end{center}
\end{document}